\chapter{Open Problems within MOS}\label{ch:open}

\begin{workplan}{\Recast}
\wpgoal{Re-rank the open problems for V2 and add the ones the pivot created.
Retitle: ``within MOS'' names the retired system.}

\wpsource{This chapter's existing text;
\texttt{TATIANA\_V1\_IDENTITY\_MAP.md}~\S10 (the ranked open decisions);
\texttt{MATH\_TOOLBOX.md}~II.1 (Tier~C material and why it was inert), III.5
(the contextuality bridge);
\texttt{TATIANA-V2/MATH/phase0\_invariant\_dimension.md}~\S4.}

\wpsteps{
  \item \textbf{Rank, and say what the ranking means:} an open problem here is
    one where a construction or a measurement is missing, not merely one where
    work remains.
  \item \textbf{\#1 --- the source of a non-exact $\eta$ for a single kernel.}
    V0 got its measured $1$-cochain from organs disagreeing. V2 has no organs.
    Without a source, Proposition~8.2 leaves growth with no address and the
    kernel cannot grow at all. This is the most important open problem in the
    book. Candidate to develop and test: disagreement between the kernel's own
    competing candidate transformations.
  \item \textbf{\#2 --- co-allocation versus separation.} When a harmonic
    residual names a cycle, does the kernel cone it (integrate) or keep the
    memories distinct? The allocation literature says excitability history
    decides, which in this formalism means it depends on the coupling weights
    --- and those weights are \emph{computed and never read}. That is a live
    bug, not merely a gap.
  \item \textbf{\#3 --- can molding erode the obstruction growth depends on?}
    Raised by the Phase~0 experiment, which bounds molding to one nudge per
    tick and so never tests it. If slow molding can dissolve a harmonic class,
    the layer handoff is a race rather than a handoff. State the experiment
    that settles it.
  \item \textbf{\#4 --- does $\mathrm{Ext}^2$ vanish for the incidence algebra
    of the $2$-complex?} This decides whether ``triangles make molding
    obstructible'' is a theorem or an analogy. Quiver categories are
    hereditary; poset categories with relations should have higher global
    dimension. Currently unverified --- do not build on it either way.
  \item \textbf{\#5 --- the MDL exchange rate}, and \textbf{\#6 --- the stall
    tolerance}: the design's only two admitted free parameters. Keep them
    listed so that a third would be conspicuous.
  \item \textbf{\#7 --- reproduction viability}
    (Chapter~\ref{ch:reproduction}), stated as currently carrying no claim.
  \item \textbf{Revisit the shelved homological material on purpose.}
    Characteristic cycles, $b$-functions and Hochschild cohomology were shelved
    as inert for one stated reason: every concept the V0 engine produced was
    simple, so the machinery had nothing to act on. A coned concept is a limit
    of a diagram and is not obviously simple, so the premise may no longer
    hold. Re-examine, carrying the standing corrections so they are not
    re-derived wrongly --- extension obstructions live in $\mathrm{HH}^3$ and
    not $\mathrm{HH}^2$; Ore extensions are unobstructed given their data; and
    a classification theorem used here needs an algebraically closed field and
    was once silently applied over $\R$.
  \item \textbf{Keep the contextuality bridge as a well-posed problem.} Both
    halves are standard separately; the intersection is not an established
    area. State the real gap sharply: a sheaf of vector spaces always admits
    the zero section, so this obstruction is about \emph{dimension}, while the
    contextuality one is about \emph{existence}.
}

\wpdone{The problems are ranked by what blocks what, and \#1 is unmistakably
identified as the one that gates the architecture.}
\end{workplan}

These are stated as problems, not as directions, because each one has a
definite answer that is currently unknown and whose resolution would
change what the system can claim. They are ordered by how much else
depends on them.

\begin{openproblem}[Non-identity restriction maps]\label{op:restrictions}
Find a procedure that learns restriction maps
$\rest{v\triang e}$ from organ data such that
(i) the maps are not all the identity,
(ii) functoriality \eqref{eq:functoriality} holds on every filled
simplex, and
(iii) the learning is closed-form or otherwise costs no gradient
training.
Closed-form Procrustes alignment satisfies (i) and (iii) but is not known
to satisfy (ii). Everything expressive about the sheaf is downstream of
this problem: while the maps are the identity, $\ker\Lap$ is the
diagonal and the sheaf is a graph Laplacian in disguise
(Proposition~\ref{prop:b0kernel}).
\end{openproblem}

\begin{openproblem}[A non-vacuous $\Hone$]\label{op:H1}
Under identity restrictions the cohomological hierarchy carries no
information: the $1$-cochain the engine emits is $\cob x$ for a global
state $x$, and a coboundary is exact by definition, so its class in
$\Hone$ vanishes identically --- for every state, every sheaf, and every
choice of restriction map. Find a source of $1$-cochains that is
\emph{not} of the form $\cob x$ --- pairwise judgements formed
independently of any global state are the candidate --- and determine
whether the resulting classes are stable enough to serve as a diagnostic.
The audit records that one proposed repair is half circular; the
circularity must be identified precisely before a fix is adopted.
\end{openproblem}

\begin{openproblem}[Does the architecture actually compound?]\label{op:slope}
The central bet is on the \emph{slope}, not the intercept: that a
self-checking, self-restructuring organiser improves with use even though
its per-inference ceiling is fixed by the underlying language model.
This is an empirical claim and it has not been tested. Design an
experiment with a pre-registered metric, a control (the same model with
no architecture), and a horizon long enough for accumulation to show, and
state in advance what result would falsify the bet.
\end{openproblem}

\begin{openproblem}[The right normaliser]\label{op:normaliser}
The Anderson--Morley bound $\AM$ over-bounds $\lambda_{\max}$, which
compresses $\rhoscore$ toward zero and contributes to its dead dynamic
range. Determine whether a tighter, still eigensolver-free bound exists
for the sheaf Laplacian with non-identity restrictions, or prove that no
degree-based bound can be tight in this regime.
\end{openproblem}

\begin{openproblem}[Stability of the diagnostics under growth]\label{op:stability}
The growth operators change the topology, and both the calibration
window and the kernel projection are topology-dependent. Give a
stability estimate: if $\opbind$ or $\opcollapse$ changes the complex by
one simplex, how much can $\rhotil$ move for a fixed state? Without such
an estimate, a change in $\rhotil$ across a growth step cannot be
attributed to the state rather than to the rewiring.
\end{openproblem}

\begin{openproblem}[Functoriality repair by hybridisation]\label{op:hybrid}
When restriction maps are obtained as projections onto an interface,
functoriality fails predictably. The canonical repair is to promote the
interface to a cell of its own, which is what hybridisation methods do in
the finite-element setting. Determine whether the same repair applies to
the cognitive complex, and what it costs: does promoting every offending
interface preserve the two-level stratification, or does it force a third
stratum?
\end{openproblem}

\begin{openproblem}[Composition factors as memory atoms]\label{op:jordanholder}
A later development formalises memory atoms as Jordan--H\"older
composition factors, with retrieval ordered by stratum, then integer
content, then structure. Determine whether the composition series is
canonical for the objects the system actually produces --- the audit
records that the associated $K_0$ classes are trivial for every such
object, which if correct means the invariant distinguishes nothing.
\end{openproblem}
